% !TEX root = ../main.tex
% !TEX program = lualatex
\documentclass[../main.tex]{subfiles}
\IfSubfilesClassLoaded{\externaldocument{\subfix{../build/main}}}{}
% =============================================================================
% File: 24_Battle_Damage_Assessment_Repair.tex
% =============================================================================
\begin{document}
\IfSubfilesClassLoaded{\chapter{EDO Study Guide}}{}
\section{Battle Damage Assessment and Repair}

\subsection{Summary}
\begin{description}
  \item[\textbf{Why it matters.}] \ac{bdar} familiarization is critical for \ac{edo} career development and wartime readiness~\autocite{edo-4-1-6-bdar-2025}.
  \item[\textbf{Repair continuum.}] Phases span damage control, afloat salvage, damage assessment, expeditionary repair, and repair at dedicated facilities~\autocite{edo-4-1-6-bdar-2025}.
  \item[\textbf{Key enablers.}] \ac{msc}, \ac{mdsu}, and \ac{supsalv} enable salvage; \ac{rmc} and \ac{nsy} teams support assessment; \ac{fdsrt}, \ac{emrf}, and \ac{sib} enable expeditionary repair~\autocite{edo-4-1-6-bdar-2025}.
  \item[\textbf{Command and control.}] \ac{c2} architecture aligns supported and supporting commands to move, task, and integrate repair forces~\autocite{edo-4-1-6-bdar-2025}.
  \item[\textbf{Contracting and funding.}] Contingency acquisition flexibilities and emergent maintenance funding enable rapid response, with fleet commanders responsible for budgeting and salvage~\autocite{edo-4-1-6-bdar-2025}.
\end{description}

\subsection{Need for BDA/R}
\begin{itemize}
  \item Renewed great power competition with the \ac{prc} and Russia, plus \ac{dprk} provocations, increases the likelihood of battle damage and contested logistics~\autocite{edo-4-1-6-bdar-2025}.
  \item \ac{bdar} planning supports readiness for high-end conflict and distributed operations~\autocite{edo-4-1-6-bdar-2025}.
\end{itemize}

\subsection{Guiding Documents and WARP}
\begin{description}
  \item[\textbf{Guidance.}] Strategic and operational guidance includes national security and defense strategy, \ac{secdef} sustainment frameworks, \ac{secnav} policy, \ac{cno} maritime sustainment strategy, and fleet planning orders~\autocite{edo-4-1-6-bdar-2025}.
  \item[\textbf{WARP.}] The Wartime Acquisition Response Plan calls on every \ac{syscom} and \ac{peo} to establish capacity, capability, data-driven actions, contingency execution, and workforce training for wartime pivot~\autocite{edo-4-1-6-bdar-2025}.
\end{description}

\subsection{BDA/R Continuum}
\begin{description}
  \item[\textbf{Damage control.}] Ship's Force combats fires, flooding, and casualties to keep the ship operable until disposition is determined~\autocite{edo-4-1-6-bdar-2025}.
  \item[\textbf{Afloat salvage.}] Rescue and assistance, towing, emergency repair, and cleanup are provided by \ac{msc}, \ac{mdsu}, and \ac{supsalv} resources~\autocite{edo-4-1-6-bdar-2025}.
  \item[\textbf{Damage assessment.}] Ship's Force, salvage forces, and ashore teams (\ac{rmc}, \ac{nsy}, \ac{navsea}) assess integrity and capability to recommend disposition~\autocite{edo-4-1-6-bdar-2025}.
  \item[\textbf{Expeditionary repair.}] Forward teams such as \ac{fdsrt} execute in-theater repairs to restore partial or full mission capability without returning to a depot~\autocite{edo-4-1-6-bdar-2025}.
  \item[\textbf{Repair.}] Major repairs occur at public or private shipyards under \ac{nsa} oversight, with options for complete repair, partial repair, or dismantling for parts~\autocite{edo-4-1-6-bdar-2025}.
\end{description}

\subsection{Disposition Factors}
\begin{itemize}
  \item Threat level and danger to rescue and repair personnel~\autocite{edo-4-1-6-bdar-2025}.
  \item Operational timelines and the minimum capabilities required to return to the fight~\autocite{edo-4-1-6-bdar-2025}.
  \item Extent of damage, mission-critical impacts, and availability of nearby facilities, escorts, and resources~\autocite{edo-4-1-6-bdar-2025}.
\end{itemize}

\subsection{Maritime Support Plan Force Packages}
\begin{description}
  \item[\textbf{MSP.}] \ac{msp} directs \ac{navsea} to provide maintenance and repair force packages across the spectrum of conflict and serve as lead integrator~\autocite{edo-4-1-6-bdar-2025}.
  \item[\textbf{Force packages.}] Afloat augmentation, afloat independent, ashore robust, ashore lean, and ashore austere packages combine personnel and equipment to meet repair requirements~\autocite{edo-4-1-6-bdar-2025}.
\end{description}

\subsection{Expeditionary Repair Capabilities}
\begin{description}
  \item[\textbf{FDSRT.}] Provides expeditionary repair, \ac{bdar} teams, contract management oversight, and fleet technical assist using broad-skill teams sourced from shipyards, \acp{rmc}, and warfare centers~\autocite{edo-4-1-6-bdar-2025}.
  \item[\textbf{SIB.}] Tiered \ac{sib} containers provide tools and materiel for afloat salvage and damage assessment (Tier 1), expeditionary repair (Tier 2), and major expeditionary repair (Tier 3)~\autocite{edo-4-1-6-bdar-2025}.
  \item[\textbf{EMRF.}] Mobile expeditionary repair facilities provide containerized capabilities for hull, structural, piping, and combat system repairs~\autocite{edo-4-1-6-bdar-2025}.
\end{description}

\subsection{Command and Control}
\begin{description}
  \item[\textbf{Supported command.}] Fleet or task force commanders provide movement, lay-down, and tasking for repair forces~\autocite{edo-4-1-6-bdar-2025}.
  \item[\textbf{Supporting command.}] \ac{navsea} provides force packages and maintenance capabilities and retains administrative control while operational tasking flows through the supported commander~\autocite{edo-4-1-6-bdar-2025}.
  \item[\textbf{Exercises.}] Ship wartime repair and maintenance exercises, \ac{bdar} exercises, expeditionary repair availabilities, salvage exercises, and voyage repair events rehearse repair integration~\autocite{edo-4-1-6-bdar-2025}.
\end{description}

\subsection{Contracting and Funding}
\begin{description}
  \item[\textbf{Contracting.}] Contingency contracting uses \ac{far}/\ac{dfars} emergency flexibilities; letter contracts authorize immediate performance as \acp{uca} when definitive contracts cannot be negotiated in time~\autocite{edo-4-1-6-bdar-2025}.
  \item[\textbf{Ship repair restrictions.}] U.S.-based ships must be repaired by U.S. companies, with foreign voyage repair limited to restoring return capability; forward deployed naval forces are the exception~\autocite{edo-4-1-6-bdar-2025}.
  \item[\textbf{Funding.}] Fleet commanders budget and execute maintenance and modernization, fund emergent repairs under restricted or technical availability programs with \ac{omn}, and finance salvage or recovery for assets under their cognizance~\autocite{edo-4-1-6-bdar-2025}.
\end{description}

\subsection{Quick Review}
\begin{itemize}
  \item \ac{bdar} spans damage control, salvage, assessment, expeditionary repair, and depot repair~\autocite{edo-4-1-6-bdar-2025}.
  \item \ac{msp} force packages and \ac{fdsrt}/\ac{emrf} capabilities enable forward repair options~\autocite{edo-4-1-6-bdar-2025}.
  \item Contingency contracting and emergent funding sources enable rapid repair action in wartime~\autocite{edo-4-1-6-bdar-2025}.
\end{itemize}

\IfSubfilesClassLoaded{
  \RenewDocumentCommand{\entryname}{}{\textbf{\color{Modern} Acronym}}
  \RenewDocumentCommand{\descriptionname}{}{\textbf{\color{Modern} Definition}}
  \printnoidxglossary[
    type=\acronymtype,
    title=Acronyms,
    style=long-booktabs
  ]}{}
\subsection{Coursebook cue: EDO role in BDAR}

For BDAR, the EDO role is to translate damage information into safe, executable
technical options under time pressure.  That requires understanding the ship's
mission priority, residual material condition, casualty-control boundaries,
technical authority, available parts and skills, and the commander's acceptable
risk.  The right answer is rarely ``restore everything''; it is the repair or
workaround that preserves mission, safety, watertight integrity, combat-system
effectiveness, mobility, and survivability with the resources available
\autocite{edo-4-1-6-bdar-2025}.

\subsection{Peacetime and wartime technical authority}

In peacetime, technical authority normally moves through the formal engineering
and certification chain before a configuration change.  In wartime or contested
conditions, the same engineering principles still matter, but the decision may
be compressed into rapid risk advice to the operational commander.  The board
answer should make the trade explicit: what is known, what is unknown, what
limit is being accepted, who has authority, and what follow-on repair or
certification is required
\autocite{edo-4-1-6-bdar-2025}.

\subsection{Contested logistics and sustainment system}

BDAR depends on people, parts, and process.  People provide damage assessment,
engineering judgment, and skilled repair.  Parts include onboard allowance,
cross-deck material, expeditionary stocks, additive or improvised repair
options, and supply-chain reachback.  Process includes reporting, triage,
technical reachback, configuration control, work authorization, and operational
risk acceptance.  Contested logistics changes the answer because parts,
connectivity, transportation, and depot access may not be available when needed
\autocite{edo-4-1-6-bdar-2025}.

\paragraph{BDAR course-of-action checklist.}
State the mission need, safety boundary, casualty data, technical authority,
parts and skills available, temporary repair limits, operational risk, and
follow-on permanent repair.

\end{document}
